\section{Hilbert tensors and minimal quotients}\label{sec:hilbert}

The purpose of this section is to turn an infinite quotient with no
infinite linear image into positivity on a base of minimal dimension.
The auxiliary compact manifold in the following statement allows the
minimality condition to be imposed before passing to an orbifold base.
For a homomorphism \(\rho\) on the fundamental group of a compact
K\"ahler manifold \(V\), write
\(\operatorname{gdim}(V,\rho)\) for the dimension of its
\(\Gamma\)-reduction.

\begin{theorem}\label{thm:hilbert-positivity}
Let \(Y\) be a connected smooth compact K\"ahler manifold of dimension
\(\ell>0\), and let
\[
 \Delta=\sum_i(1-m_i^{-1})D_i,\qquad m_i\in\N,\quad m_i\ge2,
\]
have simple normal crossing support. Denote its root orbifold by
\(\cY\). Suppose that
\[
 q:\pi_1^\orb(\cY)\longrightarrow Q
\]
is surjective, that \(Q\) is infinite, and that every finite-dimensional
complex linear representation of \(Q\) has finite image.

Suppose that a connected smooth compact K\"ahler manifold \(S\) admits
a dominant meromorphic map \(a:S\dashrightarrow Y\) and a homomorphism
\(\nu:\pi_1(S)\to Q\) with finite-index image. Assume that \(a\) and
\(\nu\) are compatible with \(q\) on a dense analytic Zariski-open
subset where \(a\) is holomorphic and takes values in
\(Y\setminus\supp\Delta\).

Finally, suppose that for every finite-index subgroup \(Q_0\le Q\)
there is a connected finite \'etale cover \(S'\to S\) such that, writing
\(\nu'\) for the induced homomorphism and \(Q'=\im\nu'\), one has
\(Q'\subset Q_0\) and
\[
 \operatorname{gdim}(S',\psi\circ\nu')\ge\ell
 \quad\text{whenever }\psi:Q'\twoheadrightarrow H,\quad
 H\text{ is infinite}.
 \tag{M}\label{hi:minimality}
\]
Then \(K_Y+\Delta\) is big.
\end{theorem}

In \eqref{hi:minimality}, the cover may depend on \(Q_0\); after it has
been chosen, the lower bound is required for every infinite quotient
of its image. Each \(Q'\) is automatically of finite index in \(Q\).
No linearity assumption is imposed on the groups \(H\).

\subsection{Holomorphic forms with Hilbert coefficients}

We use flat unitary Hilbert bundles in the usual local sense: their
transition functions are constant unitary operators. Differential
operators act on the finite-dimensional form indices and on the
Hilbert coefficients componentwise. The local elliptic and Sobolev
estimates used in Section~\ref{sec:spectral} extend to these
coefficients, as explained below.

\begin{lemma}\label{hi:finite-images}
Every finite-index subgroup \(Q_0\) of \(Q\) has only finite
finite-dimensional complex linear images. In particular,
\(H^1(Q_0,\C)=0\).
\end{lemma}

\begin{proof}
Inducing a finite-dimensional representation of \(Q_0\) to \(Q\)
produces a finite-dimensional representation of \(Q\). Its restriction
to \(Q_0\) contains the original representation, whose image is
therefore finite. The group \(Q\), and hence \(Q_0\), is finitely
generated. A nonzero homomorphism \(Q_0\to(\C,+)\) would give an
infinite unipotent two-dimensional representation, proving the last
assertion.
\end{proof}

\begin{proposition}\label{hi:coefficients}
There exist an integer \(p>0\), a separable unitary representation
\(\rho:Q\to U(\mathcal H)\) with no nonzero finite-dimensional
invariant subspace, and a nonzero holomorphic section
\[
 \eta\in H^0\bigl(\cY,\Omega_{\cY}^p\otimes\mathcal H_\rho\bigr).
\]
\end{proposition}

\begin{proof}
We treat the nonamenable and amenable cases separately.

\medskip\noindent\emph{The nonamenable case.}
Let \(U\to\cY\) be the normal orbifold covering associated with
\(\ker q\). By Theorem~\ref{thm:zero-spectrum}, in some fixed degree
\((p,0)\) there are unit vectors with spectral support in
\([0,j^{-1}]\), for \(j\to\infty\). Folding forms from \(U\) to the
compact base identifies them with sections \(s_j\) having coefficients
in the regular representation \(\lambda_Q\) on \(\ell^2(Q)\).
This is the regular-coefficient bundle interpretation of covering-space
analysis \cite[Section~6.1]{Atiyah1976}. To include the orbifold
normalization, work on a chart \([V/G]\) and put
\(H=\ker(G\to Q)\). The cover has components \([V/H]\) indexed by
the cosets of \(\im(G\to Q)\) in \(Q\). A \(G\)-invariant form with
regular coefficients is exactly the corresponding family of scalar
forms on these components. In the norm, the \(|G/H|\) regular
coordinates cancel the downstairs factor \(1/|G|\), leaving the
upstairs factor \(1/|H|\). The differential operators agree
componentwise. Thus the identification preserves the \(L^2\) norms
and operators even when \(U\) has nontrivial isotropy.

Here is the compactness statement needed for these coefficients.
For every integer \(k\ge0\), elliptic estimates give
\[
 \|s_j\|_{H^{2k}}
 \le C_k\bigl(\|(\Delta'')^k s_j\|_2+\|s_j\|_2\bigr)
 \le 2C_k.
\]
The constants do not depend on the coefficient dimension. Indeed, in
a flat chart the coefficients of the operator are finite matrices
tensored with the identity on the Hilbert space. Summing the scalar
elliptic estimates over an orthonormal expansion gives the Hilbert
estimate. The same conclusion holds on invariant orbifold charts,
whose local group orders are uniformly bounded. A finite atlas then
gives the displayed global estimate. Hilbert-valued Sobolev embedding,
obtained from Fourier inversion and Cauchy--Schwarz, bounds every
derivative uniformly on smaller charts.

Fix a nonprincipal ultrafilter. Form the Hilbert ultrapower
\(\mathcal H^\omega\) and define the values and derivatives of \(s\)
by the corresponding pointwise classes of those of \(s_j\).
Uniform bounds for the next derivative give uniform Taylor
remainders. Thus these classes are the derivatives of a smooth
section, and the exact transition identities make it a section of
the ultrapower flat bundle.

This section retains its norm. The functions
\(x\mapsto\|s_j(x)\|^2\) are uniformly bounded and uniformly Lipschitz
on the finite atlas. A finite \(\varepsilon\)-net shows that their
pointwise ultralimit is approached uniformly along the ultrafilter:
first control the finitely many net values, then use the common
Lipschitz bound. Consequently
\[
 \|s\|_2^2=\lim_\omega\|s_j\|_2^2=1.
\]
Applying the same elliptic estimates to \(\Delta''s_j\) shows
convergence to zero in every fixed \(C^r\) norm, since
\(\|(\Delta'')^{k+1}s_j\|_2\le j^{-k-1}\).
Hence \(\Delta''s=0\). Integration by parts on the compact base,
in antiholomorphic degree zero, gives \(\dbar s=0\).

The ultrapower representation has no nonzero finite-dimensional
invariant subspace. Otherwise such a subspace would have finite
image, by the hypothesis on \(Q\), and would be fixed by a
finite-index subgroup \(Q_0\). This subgroup is nonamenable.
For a finite generating set \(F\subset Q_0\) and some
\(\varepsilon>0\),
\[
 \sum_{g\in F}\|\lambda(g)v-v\|^2
 \ge\varepsilon\|v\|^2
\]
on its regular representation. Indeed, failure would give unit
almost invariant vectors; their squared absolute values would be
almost invariant probability measures in \(\ell^1(Q_0)\), yielding
an invariant mean. Summing over copies gives the same
inequality on the restriction of \(\lambda_Q\) to \(Q_0\).
A nonzero \(Q_0\)-fixed ultrapower vector, represented by a bounded
sequence of limiting norm one, contradicts this inequality after
taking the ultralimit. There are only finitely many terms on its
left side.

Finally, restrict to the closed span of all coefficient values of
the section and their \(Q\)-translates. Countable charts, continuity,
and countability of \(Q\) make this subspace separable. It still
contains the nonzero section and has no finite-dimensional invariant
subspace.

\medskip\noindent\emph{The amenable case.}
An infinite amenable group does not have property~\((T)\): its regular
representation has almost invariant vectors and no invariant vector.
Shalom's reduced-cohomology theorem therefore gives a unitary
representation with nonzero reduced first cohomology
\cite[Theorem~4.2]{Shalom2006}.
We may restrict to a separable invariant subspace generated by the
values of a cocycle on the countable group.

The closed span of the finite-dimensional invariant subspaces
contributes no reduced first cohomology. To see this, decompose it
orthogonally into finite-dimensional irreducible summands. On any
finite sum the representation has finite image. Its kernel \(Q_0\)
has \(H^1(Q_0,\C)=0\), by Lemma~\ref{hi:finite-images}; restriction of
a cocycle to that kernel is zero, and averaging over the finite
quotient makes the cocycle a coboundary. Finite partial sums
approximate a cocycle on a finite generating set, so the same
vanishing holds for reduced cohomology of the closed sum.
The orthogonal complement therefore has a cocycle \(b\) that is not
a limit of coboundaries and has no finite-dimensional subrepresentation.

Pullback by the surjection \(q\) preserves this nonvanishing:
approximating the pullback cocycle on lifts of finitely many
generators of \(Q\) would approximate \(b\) itself. The associated
flat affine Hilbert bundle on \(\cY\) has a smooth section, constructed
by local sections and a partition of unity, averaged in the finite
orbifold charts. Its covariant differential is a smooth closed
Hilbert-valued one-form \(\alpha\) representing the pulled-back
cocycle.

The harmonic projection of \(\alpha\) is nonzero. Otherwise, the
orthogonal decomposition of the closed de Rham Hilbert complex
would put the closed form \(\alpha\) in the \(L^2\) closure of
exact forms. For fixed \(t>0\), the heat operator
\(e^{-t\Delta}\) commutes with the
differential and maps \(L^2\) convergence to \(C^0\) convergence,
by the dimension-independent elliptic estimates just used.
Thus \(e^{-t\Delta}\alpha\) would be a uniform limit of smooth exact
forms. Periods along finitely many fixed generator loops, with flat
parallel transport of coefficients, would approximate its cocycle
by coboundaries. Letting \(t\downarrow0\) gives the same conclusion
for \(\alpha\), contrary to the choice of \(b\).
The orthogonal decomposition of the closed de Rham Hilbert complex
now gives a nonzero harmonic one-form; local elliptic regularity
makes it smooth.

The K\"ahler identities for flat unitary coefficients split this
harmonic form into types. A nonzero \((1,0)\)-part is holomorphic.
If only the \((0,1)\)-part is nonzero, complex conjugation gives a
holomorphic \((1,0)\)-form with conjugate coefficients. The conjugate
representation again has no finite-dimensional invariant subspace.

In either case, degree zero is impossible: a holomorphic section of
a flat unitary Hilbert bundle on a compact K\"ahler orbifold is
parallel, by integration of the Laplacian of its squared norm.
A nonzero parallel section would give an invariant line.
\end{proof}

\subsection{The line associated with a tensor}

Proposition~\ref{hi:coefficients} supplies a nonzero tensor, but
nonvanishing alone does not contradict specialness. We next isolate
its coefficient line and measure the variation of that line by a
positive curvature current. This connects the group to cotangent
positivity, as in the finite-dimensional setting of
\cite{BrunebarbeKlinglerTotaro2013}; the Hilbert coefficient space
and the subsequent minimal-rank argument require the constructions
proved here.

For a finite-index subgroup \(\Gamma\le Q\), let
\(\cY_\Gamma\to\cY\) be the cover associated with \(q^{-1}(\Gamma)\).
All finite orbifold covers considered in this subsection and the
next are of this form.

Call a holomorphic tensor
\[
 \tau\in H^0\bigl(\cY_0,
       (\Omega_{\cY_0}^1)^{\otimes m}\otimes\mathcal B_\rho\bigr)
\]
a \emph{line tensor} if the map from the dual of its cotangent
tensor factor to its Hilbert coefficient space has generic rank one.
Here \(\cY_0\) is a connected finite orbifold cover of \(\cY\),
and \(\rho\) is a unitary representation of its corresponding
finite-index subgroup of \(Q\). On the \(Q\)-cover, the coefficient
lines define a holomorphic map to \(\PP(\mathcal B)\) away from an
analytic degeneracy locus. We always replace \(\mathcal B\) by the
closed span of these lines.

\begin{lemma}\label{hi:plucker}
There is a line tensor whose coefficient lines have infinite-dimensional
closed span.
\end{lemma}

\begin{proof}
Let \(\eta\) be given by Proposition~\ref{hi:coefficients}, and let
\(r\) be the generic rank of
\((\Omega_{\cY}^p)^*\to\mathcal H\).
Taking its \(r\)-th exterior power gives a nonzero rank-one tensor
with coefficient space \(\bigwedge^r\mathcal H\).
Antisymmetrization embeds its finite-dimensional tensor factor into
\((\Omega_{\cY}^1)^{\otimes pr}\).

If the resulting decomposable lines spanned a finite-dimensional
space, \(Q\) would act on that span through a finite group. The orbit
of one of these lines would be finite. A decomposable exterior line
determines its \(r\)-dimensional original coefficient plane, so the
corresponding orbit of planes would also be finite. Their sum would
be a nonzero finite-dimensional invariant subspace of \(\mathcal H\),
a contradiction.
\end{proof}

\begin{lemma}\label{hi:saturated-line}
A line tensor determines a saturated orbifold line subsheaf
\[
 \cL\subset(\Omega_{\cY_0}^1)^{\otimes m}
\]
and a positive curvature current on \(\cL\). On the regular locus
this current is the pullback of the Fubini--Study form of its
coefficient-line map.
\end{lemma}

\begin{proof}
Work in a uniformizing chart with a flat coefficient trivialization.
Any nonzero scalar coefficient of \(\tau\) determines its
meromorphic direction in the finite-rank tensor bundle. All scalar
coefficients have this direction, since the coefficient rank is one.
These local meromorphic directions agree on overlaps. Their
saturations are coherent rank-one subsheaves of the tensor bundle.
A saturated subsheaf of a locally free sheaf is reflexive in rank
one, and on a smooth chart a rank-one reflexive sheaf is a line
bundle. The invariant constructions give the asserted orbifold line.

Write a local line generator as \(e\). Off its codimension-two
degeneracy set the tensor has the form \(e\otimes u\), with \(u\)
Hilbert-valued and holomorphic. There are no divisorial poles, by
saturation. The Hilbert-valued Hartogs extension, obtained from the
Cauchy integral formula on transverse polydiscs, extends \(u\)
across the remaining set.

The weights \(\log\|u\|^2\) are plurisubharmonic. If
\(e_\beta=g_{\alpha\beta}e_\alpha\), then
\(u_\beta=g_{\alpha\beta}^{-1}u_\alpha\), so
\[
 \log\|u_\beta\|^2
 =\log\|u_\alpha\|^2-\log|g_{\alpha\beta}|^2.
\]
They therefore define the metric
\(\|e_\alpha\|^2=\|u_\alpha\|^{-2}\) on \(\cL\), with positive
curvature. Where \(u\ne0\), its curvature is
\(\dd\dc\log\|u\|^2\), the pulled-back Fubini--Study form.
A positive power of \(\cL\) kills the finitely many chart
stabilizer actions and descends to an ordinary line bundle on the
coarse space; its weights descend as psh weights.
\end{proof}

The degeneracy loci used here and below are analytic even though the
coefficient space may be infinite-dimensional. Locally their rank
conditions are the vanishing of holomorphic scalar minors.
The ideal of their germs at a point is generated by finitely many
of those minors. On sufficiently small representatives of the
finitely many local components, every remaining minor vanishes by
the identity theorem. This proves local analyticity of the common
zero set; it does not require a global finite list of coefficients.

\subsection{Minimal rank and discrete monodromy}

Among all line tensors with infinite-dimensional span on finite
orbifold covers, choose one for which the generic complex rank \(s\)
of the coefficient-line map is minimal. This is a minimum in the
finite set \(\{1,\ldots,\ell\}\): rank zero would make the map
constant on its connected covering, giving a one-dimensional span.
Let \(\Gamma\le Q\) be the subgroup defining the selected cover,
let \(\mathcal B\) be its separable coefficient span, and write
\[
 F:U^\circ\longrightarrow\PP(\mathcal B)
\]
on the connected regular open subset of the \(Q\)-cover. We may
delete the chart isotropy locus as well. These deletions are proper
analytic sets and do not disconnect the covering.

\begin{proposition}\label{hi:discrete-monodromy}
After replacing \(\Gamma\) by a finite-index subgroup, the action on
the regular image germs of \(F\) has infinite discrete image in a
Lie group acting properly on a connected complex manifold of
dimension \(s\).
\end{proposition}

\begin{proof}
\medskip\noindent\emph{The manifold of image germs.}
At each point of \(U^\circ\), the holomorphic constant-rank theorem
writes \(F\) as a submersion onto an embedded \(s\)-dimensional
plaque. This theorem also holds for the present Hilbert target:
choose a finite-dimensional projection of rank \(s\), use it for
source coordinates, and observe that all derivatives in the
remaining source directions vanish.

Let \(Z\) consist of the image points together with the germs of
these embedded plaques at those points. A plaque supplies its own
coordinate chart on \(Z\). The map \(U^\circ\to Z\) is an open
holomorphic submersion, so \(Z\) is connected and second countable.
It is Hausdorff. Indeed, two germs over the same image point can
both be written as graphs over one finite-dimensional transverse
projection and one small connected ball. If their plaque
neighborhoods shared a germ, the graph functions would agree on an
open set and hence everywhere on that ball. The original germs
would be equal. Points lying over different image points are
separated by the Hilbert projective space.

The natural immersion \(\iota:Z\to\PP(\mathcal B)\) gives \(Z\)
a K\"ahler metric. We use its intrinsic Riemannian length distance,
which induces the manifold topology and is locally compact even when
the metric is incomplete. Its isometry group acts properly
\cite[Corollary~4]{Manoussos2010} and is a Lie group
\cite[author's version, Theorem~6.3]{MatveevTroyanov2017}.
Let \(I\) be the closure of the \(\Gamma\)-image in this isometry
group with its compact-open topology. Its elements are holomorphic.
Indeed, if holomorphic isometries \(g_j\) converge to \(g\), choose
a relatively compact neighborhood whose image under \(g\) lies
compactly inside one target complex chart. Eventually the \(g_j\)
also take its closure into that chart, and their coordinate maps
converge locally uniformly. The Weierstrass theorem makes \(g\)
holomorphic there. Inversion is continuous in the isometry group,
so the same argument applies to \(g^{-1}\).

\medskip\noindent\emph{A locally compact unitary extension.}
Let \(J\) be the closure of the simultaneous image of \(\Gamma\) in
\(I\times U(\mathcal B)\), with the strong group topology on the
unitary factor. Every \((g,A)\in J\) satisfies
\(A\iota(z)=\iota(gz)\) projectively, by continuity.
The projection \(J\to I\) is proper and surjective.
For properness, suppose that the isometry components of a sequence
converge. Choose countably many unit vectors in coefficient lines
whose span is dense in \(\mathcal B\). The images of each such
vector lie in convergent projective lines. Compactness of the
scalar circle gives a convergent subsequence of the unit vectors.
Diagonal extraction, performed also for inverse operators, gives
strong convergence to a unitary operator. This proves properness,
since the spaces involved are metrizable. The image is consequently
closed and contains the dense \(\Gamma\)-image, proving surjectivity.

The kernel is a closed subgroup of the scalar circle. A kernel
element preserves every coefficient line; a unitary operator has
the same eigenvalue on two nonorthogonal eigenlines. Nearby lines
are nonorthogonal, so connectedness of \(Z\) makes this scalar the
same everywhere. Their span is all of \(\mathcal B\).
It follows that \(J\) is locally compact. It has no small subgroups:
use such a neighborhood in the Lie quotient \(I\), then one in the
scalar kernel. The locally compact no-small-subgroups theorem
therefore makes \(J\) a Lie group
\cite[author's version, Corollary~1.5.8]{Tao2014}.

Its adjoint representation, complexified, has finite image on
\(\Gamma\), by Lemma~\ref{hi:finite-images}. Pass to its kernel in
\(\Gamma\), and replace both \(I\) and \(J\) by the corresponding
closures. They are open finite-index subgroups of the old closures,
with unchanged identity components; \(J\to I\) remains surjective
and proper. The \(Q\)-cover and the map \(F\) are unchanged.
Every element of the new \(J\) centralizes \(J^\circ\), since
its conjugation has identity differential on the connected group.
In particular, \(J^\circ\) is abelian.

\medskip\noindent\emph{Spectral projection of holomorphic data.}
The spectral theorem for the continuous unitary representation of
the connected abelian Lie group \(J^\circ\), in its commutative
\(C^*\)-algebra formulation \cite{Williams2007}, gives a direct integral
\[
 \mathcal B=\int^\oplus \mathcal B_\chi\,\dd\mu(\chi),\qquad
 t|_{\mathcal B_\chi}=\chi(t)\Id .
\]
We take \(\mu\) finite, as is possible because \(\mathcal B\) is separable.
Because \(\Gamma\) centralizes \(J^\circ\), its operators are
decomposable in this integral \cite{Williams2007}.
Write \(j_\gamma\) for the simultaneous image of \(\gamma\) in \(J\),
and put
\[
 D=\{j_\gamma:\gamma\in\Gamma\}\cap J^\circ.
\]
Countability of \(\Gamma\) allows the group identities to hold on
one common full-measure set, giving unitary representations
\(\rho_\chi\) of \(\Gamma\). On that set we also impose the countably
many identities
\[
 \rho_\chi(\gamma)=\chi(j_\gamma)\Id
 \qquad\text{when }j_\gamma\in D.
\]

We explain carefully how the holomorphic tensors pass to these
fibres. These are almost-everywhere decompositions of a countable
family of Taylor coefficients; no bounded point-evaluation projection
from the direct integral is asserted. Let
\(\varpi:U^\circ\to Z\) be the submersion. Choose a countable
uniformizing atlas for \(\cY_\Gamma\) and all its lifts to the
\(Q\)-cover, still a countable family. Choose also local nonvanishing
holomorphic vector representatives \(v_\alpha\) of \(\iota\) on a
countable cover of \(Z\). On a source polydisc mapping
into the domain of \(v_\alpha\), the finitely many tensor coefficients
satisfy vector identities
\[
 \tau_a(x)=c_a(x)v_\alpha(\varpi(x)),
\]
with holomorphic scalar functions \(c_a\), obtained after shrinking
by applying a bounded functional nonzero on \(v_\alpha\).
On suitable overlap
domains the equivariance identities have the form
\[
 v_\beta(\gamma z)
 =c_{\gamma,\alpha\beta}(z)\rho(\gamma)v_\alpha(z),
\]
where \(c_{\gamma,\alpha\beta}\) is a holomorphic unit.
Here \(\rho\) is the selected unitary representation of \(\Gamma\).
These identities are kept as vector identities before projectivizing.

Expand the tensor coefficients and the \(v_\alpha\) in
Hilbert-valued power series. On every strictly smaller polydisc the
sums of coefficient norms times the monomial radii are finite.
Cauchy--Schwarz in the spectral variable, followed by Tonelli,
gives the corresponding absolute sums in
\(\mathcal B_\chi\) for almost every \(\chi\). The projected series
are therefore holomorphic there. After restricting to countably
many smaller overlap polydiscs, normal convergence permits
composition with the fixed coordinate maps and multiplication by
the scalar functions above term by term. For each resulting coefficient
sum, its partial sums converge in the direct-integral norm, while
its fibre series converges absolutely almost everywhere. A subsequence
converging almost everywhere identifies the two limits. There are
only countably many such sums and Taylor-coefficient identities,
including the tensor transition and orbifold isotropy identities.
They therefore hold, together with the group identities above,
on one common full-measure set and give global
holomorphic tensors on
\(\cY_\Gamma\) of coefficient rank at most one, and holomorphic
projective maps
\[
 F_\chi:Z\setminus E_\chi\longrightarrow\PP(\mathcal B_\chi)
\]
of rank at most \(s\), where \(E_\chi\) is the common zero locus of
the projected representatives.

The countable family of Taylor coefficients of the tensor on the
lifted charts has closed span \(\mathcal B\): the series give all
coefficient values, while Cauchy integrals express their coefficients
as limits of linear combinations of those values. The orthogonal
projection onto the fibrewise complement of the projected coefficient
spans is a measurable field, obtained by countable Gram--Schmidt
orthogonalization. Its direct integral annihilates this dense family
and is therefore zero. Thus the projected Taylor coefficients span
\(\mathcal B_\chi\) almost everywhere. By the same series and Cauchy
formulas, so do the projected coefficient values. On every nonzero
such spectral fibre the projected tensor is therefore nonzero,
has generic coefficient rank one, and is a line tensor whose
coefficient lines span \(\mathcal B_\chi\). We henceforth restrict
to these nonzero fibres.

\medskip\noindent\emph{The identity component cannot move an image germ.}
The subgroup \(D\) is dense in \(J^\circ\), since that component is open.
It acts by the character \(\chi\) on each spectral fibre, so
\(F_\chi\) is invariant under its action on \(Z\).
The closed analytic set \(E_\chi\) is invariant as well.
Continuity extends both invariances to \(J^\circ\).

Suppose that \(J^\circ\) acts nontrivially on \(Z\). If some
\(F_\chi\) had rank \(s\) at \(z\notin E_\chi\), a finite-dimensional
target projection of full rank and the inverse function theorem
would make \(F_\chi\) locally injective at \(z\). The connected
\(J^\circ\)-orbit of \(z\) stays in the domain and in its fibre by
invariance. The point \(z\) is isolated in that fibre, and its
singleton is open there by isolation and closed because \(Z\) is Hausdorff.
The orbit is therefore \(\{z\}\). Applying this at every point of
the nonempty open full-rank locus, each element of \(J^\circ\)
fixes that open set. It is holomorphic, so the identity theorem on
connected \(Z\) makes it the identity everywhere, a contradiction.
Thus every \(F_\chi\) has rank strictly less than \(s\).
Minimality of \(s\) then forces
\(\mathcal B_\chi\) to be finite-dimensional almost everywhere.
By Lemma~\ref{hi:finite-images}, \(\rho_\chi(\Gamma)\) is finite.
The character \(\chi\) consequently has finite image on the dense
subgroup of \(J^\circ\) just used. Continuity and connectedness
force \(\chi\) to be trivial on \(J^\circ\).
Hence \(J^\circ\) acts trivially on \(\mathcal B\).
It then acts trivially on \(Z\): each orbit maps to a point under
the immersion \(\iota\), and is connected. This contradicts the
assumed nontrivial action.

We have proved that \(J^\circ\) acts trivially on \(Z\).
The surjective proper Lie-group homomorphism \(J\to I\) maps
identity component onto identity component; since \(I\) is an
effective isometry group, \(I^\circ=\{1\}\).
Thus \(I\) is discrete and acts properly.

Its \(\Gamma\)-image is infinite. Otherwise a finite-index subgroup
would fix every image germ and would act on \(\mathcal B\) by
scalars. Its scalar character has finite image; on a further finite
orbifold cover the tensor would be untwisted. An untwisted
Hilbert-valued tensor on a compact orbifold has finite-dimensional
coefficient span: its scalar tensor space is finite-dimensional,
and expansion in a basis writes the tensor as a finite sum of
scalar tensors times fixed Hilbert vectors. This contradicts our
choice of the line tensor.
\end{proof}

\subsection{Compact fibres from transverse volume}

Discrete monodromy gives a genuine quotient of the image-germ
manifold. To use minimality, its local fibres must determine compact
subvarieties with finite group image. The following argument obtains
these subvarieties from the mass of the positive current; only the
source is assumed compact.

\begin{lemma}\label{hi:compact-fibres}
Let \(V\) be a connected smooth compact K\"ahler manifold of
dimension \(d\), and let \(V^\circ\) be the complement of a proper
analytic subset. Let a discrete group \(I\) act properly and
effectively by holomorphic isometries on a connected K\"ahler
manifold \(Z\) of dimension \(s>0\).
Suppose that \(f:V^\circ\to Z/I\) has local holomorphic lifts to
\(Z\) of rank \(s\), whose continuation is compatible with a
homomorphism \(\rho:\pi_1(V)\to I\).
If a positive closed \((1,1)\)-current on \(V\) restricts to
a positive constant multiple of \(f^*\omega_Z\), then
\begin{equation}
 \operatorname{gdim}(V,\rho)\le s.
 \label{hi:gamma-upper-bound}
\end{equation}
\end{lemma}

\begin{proof}
Let \(T\) be the current and let \(\omega\) be a K\"ahler form on
\(V\). The mass of \(T^s\wedge\omega^{d-s}\) on \(V^\circ\) is
finite; this is the positive-current mass bound
\cite{Boucksom2002}. One can see the bound directly
by writing \(T=\alpha+\dd\dc\varphi\), choosing \(C\) with
\(\beta=\alpha+C\omega>0\), and truncating
\(\varphi_j=\max(\varphi,-j)\). The bounded-potential products satisfy
\[
 \int_V(\beta+\dd\dc\varphi_j)^s\wedge\omega^{d-s}
 =\int_V\beta^s\wedge\omega^{d-s}.
\]
On \(\{\varphi>-j\}\cap V^\circ\) they equal
\((T+C\omega)^s\wedge\omega^{d-s}\), which dominates
\(T^s\wedge\omega^{d-s}\). Exhaustion gives the bound.

Use the free locus of the orbifold \(Z/I\) for the coarea formula.
The discarded isotropy values, and their inverse images under the
local submersions, have measure zero. With normalized volumes,
\[
 f^*\frac{\omega_Z^s}{s!}\wedge
       \frac{\omega^{d-s}}{(d-s)!}
 =J_f\,\frac{\omega^d}{d!},
\]
where \(J_f\) is the real normal Jacobian. Its value is the product
of the squared nonzero complex singular values of the differential.
The coarea formula and the mass bound show that almost every
fibre has finite \((d-s)\)-dimensional complex volume in \(V^\circ\).
The formula is applied to a countable disjoint measurable
decomposition subordinate to target charts, so it counts whole
fibres and requires neither compactness nor bounded geometry of
the target. Moreover, this yields a full-measure set of source
points: the inverse image of the exceptional null set has integral
of \(J_f\) equal to zero, and \(J_f>0\) on the submersion locus.

For a good free value, its fibre is closed analytic and pure of
dimension \(d-s\) in \(V^\circ\). Its integration current has finite
mass. Extension by zero across \(V\setminus V^\circ\) is positive
and closed, by the extension theorem of El Mir
\cite{ElMir1984}.
Siu's analyticity theorem for positive Lelong superlevel sets
\cite{Siu1974} then makes its closure analytic of pure dimension
\(d-s\): the level set with Lelong number at least one contains
the fibre, has dimension at most \(d-s\), and agrees with the
smooth fibre on \(V^\circ\). Take its components meeting
\(V^\circ\). Equivalently, one may apply the finite-volume
closure theorem directly \cite{ElMir1984}.

Let \(\widetilde C\to C\) resolve one such closure component.
The preimage of \(C\cap V^\circ\) has proper analytic complement
in the smooth manifold \(\widetilde C\), and its inclusion induces
a surjection on fundamental groups: every loop may be perturbed
off an analytic set of real codimension at least two.
Along the open fibre the continued lift to \(Z\) is constant.
Thus its monodromy fixes one point of \(Z\), and lies in a finite
stabilizer. The same is true of the image of
\(\pi_1(\widetilde C)\), by that surjection.

The \(\Gamma\)-reduction contracts these compact subvarieties through
very general points, by Theorem~\ref{thm:gamma-reduction}.
The full-measure set just obtained meets the complement of its
countable union of proper analytic exceptional sets. Its general
fibres consequently have dimension at least \(d-s\), proving
\eqref{hi:gamma-upper-bound}.
\end{proof}

\subsection{Descent of the big line}

\begin{proof}[Proof of Theorem~\ref{thm:hilbert-positivity}]
Choose the minimal line tensor and the finite-index subgroup
\(\Gamma\) of Proposition~\ref{hi:discrete-monodromy}.
Let \(\cL\) and its positive curvature current be as in
Lemma~\ref{hi:saturated-line}. The current has rank \(s\) on the
regular open set.

Suppose that \(s<\ell\). Apply \eqref{hi:minimality} with
\(Q_0=\Gamma\), obtaining \(S'\) with image \(Q'\subset\Gamma\).
The map \(S'\dashrightarrow Y\) lifts meromorphically to the
coarse space of \(\cY_\Gamma\). Indeed, the monodromy condition
gives a lift on the complementary open set; the selected component
of the normalized finite pullback extends its graph.
Resolve this map on a smooth compact K\"ahler modification \(V\)
of \(S'\).

The descended power of \(\cL\) pulls back to an ordinary line
bundle on \(V\). Its psh weights pull back, since the map is
dominant, giving a positive current. On a regular analytic
Zariski-open subset the local coefficient-line maps give
\[
 V^\circ\longrightarrow Z/I
\]
of rank \(s\), with monodromy the composite
\(\pi_1(V)\to Q'\to I\). The pulled-back current is a positive
constant multiple of its transverse K\"ahler form.
Lemma~\ref{hi:compact-fibres} gives a \(\Gamma\)-dimension at most
\(s\). The image of \(Q'\) in \(I\) is infinite: it has finite
index in the infinite image of \(\Gamma\).
This image is an infinite quotient to which \eqref{hi:minimality}
applies. Bimeromorphic invariance of the \(\Gamma\)-reduction
therefore gives a dimension at least \(\ell\), a contradiction.
Hence \(s=\ell\).

Take a finite Galois orbifold cover
\(\widehat{\cY}\to\cY\) dominating \(\cY_\Gamma\), and pull
\(\cL\) to it. If \(G\) is its deck group, form
\[
 \mathcal M=\bigotimes_{\gamma\in G}\gamma^*\cL.
\]
The permutation action gives a canonical \(G\)-linearization.
All its factors inject into the pullback of the same cotangent
tensor bundle. Multiplication in the symmetric algebra gives a
nonzero morphism from \(\mathcal M\) to the \(|G|\)-th symmetric
power of that tensor bundle, and hence into a cotangent tensor power.
It is nonzero because a product of nonzero vectors in a symmetric algebra is
nonzero. Equivariant descent gives an orbifold line \(\mathcal M_0\)
and a nonzero cotangent-tensor morphism on \(\cY\).

The tensor product of the translated metrics is invariant.
Its curvature is positive and strictly positive on a nonempty
open set, since \(s=\ell\). A power of \(\mathcal M_0\) kills
the chart stabilizers and descends to an ordinary line bundle
\(A\) on \(Y\) with the same positivity property. The volume
criterion \cite[author's preprint, Theorem~1.2]{Boucksom2002} makes \(A\) big.
The Iitaka map of a big line bundle makes \(Y\) Moishezon.
The projectivity theorem for K\"ahler Moishezon manifolds
\cite{Moishezon1967} therefore makes \(Y\) projective.

It remains to interpret the tensor inclusion on an adapted
projective cover. Choose the smooth Galois Kawamata cover
\(\pi:Y'\to Y\) supplied by \cite[Lemma~5.2]{CampanaPaun2019}.
It is projective, being finite over \(Y\), and its ramification
order over \(D_i\) is exactly the denominator \(m_i\).
Near that component the root chart and the cover are
\(z_i=u_i^{m_i}=w_i^{m_i}\). Choosing the local root lift
\(u_i=w_i\), the cotangent generator pulls back as
\[
 \dd u_i=\dd w_i.
\]
Along an auxiliary ramification divisor \(z_j=w_j^{a_j}\), the
orbifold coefficient is zero and the generator is
\(\dd z_j=a_jw_j^{a_j-1}\dd w_j\).
These are precisely the local generators of the adapted orbifold
cotangent bundle in \cite[Definition~5.3]{CampanaPaun2019}.
Thus a suitable tensor power of the descended inclusion gives
a nonzero map
\[
 \pi^*A\longrightarrow
 \bigl(\pi^*\Omega^1(Y,\Delta)\bigr)^{\otimes M}
\]
on this projective adapted cover. The big-line criterion for orbifold
cotangent tensors \cite[Theorem~7.11]{CampanaPaun2019} now implies
that \(K_Y+\Delta\) is big. Its hypotheses are satisfied because
projectivity and the ordinary big line \(A\) have already been
established. No pseudoeffectivity assumption is needed in that
criterion.
\end{proof}
